\thispagestyle{plain}

\noindent
\begin{tcolorbox}[
  colback=nexanavy,
  colframe=nexanavy,
  arc=4pt,
  boxrule=0pt,
  left=15pt,
  right=15pt,
  top=15pt,
  bottom=15pt
]
  \begin{minipage}[t]{0.7\textwidth}
    \color{white}
    {\huge\bfseries Nexa University}\\[6pt]
    {\large\color{nexateal!40}\bfseries Office of Academic Assessment}\\[12pt]
    {\color{white}\Large Student Group Project Peer Evaluation Rubric}\\[4pt]
    {\color{nexaamber}\small A Standardized Guide and Instrument for Peer Performance Review}
  \end{minipage}%
  \begin{minipage}[t]{0.3\textwidth}
    \raggedleft
    \color{white}
    \faGraduationCap\ \textbf{EDU-R-2026}\\[4pt]
    \small
    Fall Semester 2026\\
    Version 2.4
  \end{minipage}
\end{tcolorbox}

\vspace{15pt}

\noindent
\begin{tcolorbox}[
  colback=nexabg,
  colframe=nexateal,
  arc=3pt,
  boxrule=0.8pt,
  left=10pt,
  right=10pt,
  top=8pt,
  bottom=8pt
]
  \small\color{nexacharcoal}
  \textbf{Abstract \& Purpose:} This rubric and assessment instrument are designed to facilitate fair, rigorous, and constructive peer evaluation in student group projects. By assessing teammates across five core dimensions---Contribution to Work, Quality of Output, Communication, Collaboration \& Dynamics, and Reliability---instructors can diagnose team dynamics and calculate adjusted individual grades. This document outlines the evaluation framework, the scoring rubric, and the submission form.
\end{tcolorbox}
